\section{漏洞检测为什么难：语义上下文、评测噪声与动态验证}

\subsection{内核案例：局部可疑不代表全局成立}
课程用内核代码示例解释了漏洞判定难点：某函数从包头读取长度后分配缓冲区并复制数据，若长度与真实 payload 不一致可能导致越界写。困难不在于模式识别本身，而在于\textbf{需要确认可达路径、调用方约束和跨函数数据流}。

\begin{figure}[H]
\centering
\includegraphics[width=0.88\textwidth]{frame-05-kernel-vuln.jpg}
\caption{内核漏洞讲解片段：同一函数的风险取决于全局调用上下文 \protect\footnotemark}
\end{figure}
\footnotetext{视频画面时间区间：01:06:30--01:07:10。}

\begin{importantbox}{主讲人的关键判断}
\textit{``This function on its own is not vulnerable or non-vulnerable.''}  
单函数分类往往不是良定义问题。是否漏洞成立，取决于调用方是否已约束输入、路径是否可达、运行环境是否满足触发条件。
\end{importantbox}

\subsection{漏洞数据库驱动评测的天然难题}
公开漏洞库（如 CVE 生态）看起来是理想数据源：有编号、有描述、有补丁。但将其直接转换为 ``函数级分类任务'' 会遇到严重问题：输入上下文该截到哪里、正负样本如何定义、补丁是否只对应一个根因、标签是否跨版本稳定。

\begin{warningbox}{“下载漏洞库就能训练检测器”通常过于乐观}
课堂明确指出，这是最自然但也最容易踩坑的路径。若不先定义严谨任务边界，模型可能学到项目风格、提交习惯或模板语句，而不是真正的可利用性语义。
\end{warningbox}

\subsection{动态评测的必要性}
Sutton 提到 DARPA 相关竞赛与 Meta 安全评测方向，核心共同点是从静态标签转向可执行验证。对于漏洞发现任务，动态评测更接近真实目标：输入触发、运行证据、崩溃位置、可复现性。

\begin{knowledgebox}{Fuzzing 与 Agent 的互补关系}
\begin{itemize}
    \item Fuzzing 擅长高吞吐随机探索；
    \item Agent 擅长语义引导与路径假设；
    \item 二者结合可形成 ``广覆盖 + 深路径'' 的混合搜索。
\end{itemize}
这也是讲座后段反复强调的方向：AI 不替代 fuzzing，而是把搜索推向 fuzzer 难以触达的深层逻辑分支。
\end{knowledgebox}

\subsection{从静态精度转向任务成功率}
安全场景里，最终指标不该是 ``这段代码像不像漏洞''，而应是 ``是否能构造触发输入并给出可审计证据''。这会把研究焦点从文本分类转到\textbf{可执行任务成功率}，与 Agent 范式天然一致。

\begin{importantbox}{对研究设计的直接启示}
未来安全 Agent 论文若只报告静态分类分数，很难支撑工程价值；至少应增加动态触发率、复现稳定性、误报处置成本等指标。
\end{importantbox}

\subsection{本章小结}
漏洞检测难在 ``上下文依赖 + 可执行验证''，而不是 ``模式匹配不足''。这也解释了为什么安全方向更需要 Agent 化系统，而不只是更强的代码文本模型。

